% ===========================================================================
\section{Related Work}
\label{sec:related}
% ===========================================================================

\subsection{Fuel-based and fleet-based transport inventories}

The IPCC 2006 Guidelines describe both approaches for 1A3b and are explicit that the fuel-based
Tier~1 is preferred where national fuel statistics are reliable and fleet data are not
\citep{ipcc2006energy}. That guidance is often inverted in practice, because a fleet-based estimate
produces the disaggregation that transport policy wants and a fuel-based one does not.

The position taken here is that the two are not alternatives at the same level. The fuel
figure determines the total; the fleet model determines the split. Running the fleet model
\emph{inside} the fuel constraint gives the policy-relevant disaggregation without letting it
set the magnitude, which is the arrangement the guidelines' preference implies but which is
rarely stated as such.

\subsection{Global products and their transport scope}

EDGAR \citep{edgar2025booklet,crippa2020} reports Pakistani transport across all modes---road, rail,
domestic aviation and navigation together. Comparing that against a road-only estimate is a
scope error that will make any road inventory look low; Section~\ref{sec:validation} walks
the scope component by component instead, and labels the result honestly as fiscal-year blended
and partially scope-harmonized.

Climate TRACE \citep{climatetrace} constructs road transport from asset and activity
inference. Both are global models applying globally consistent assumptions, which is their
strength across countries and their limitation within one: neither can know that Pakistan's
national diesel figure includes agricultural and captive-generation fuel, which is the single
largest scope question in this sector.

\subsection{Off-road machinery, and where it belongs}

IPCC reports off-road mobile machinery outside road transport, and not all under one
category: agricultural tractors and plant go to 1A4c\,ii, while construction and industrial
machinery belongs under 1A2 (manufacturing and construction, gvii). This is easy to state
and easy to get wrong in a country where the diesel that runs all of it comes from the same
retail pumps as the diesel that runs the trucks and is not separately reported---and this
inventory gets it only partly right: its off-road class inherits a residual bucket that
mixes both, and is reported under the mixture-declaring name \emph{Mixed off-road
(1A4c/1A2)} with the misassignment declared in the completeness table and priced as a
reporting-recode scenario rather than resolved (Section~\ref{sec:limitations}).

The usual treatments are to leave it inside road transport, which overstates 1A3b and leaves
1A4c empty, or to estimate it separately from a machinery population, which double-counts the
fuel unless the road control total is reduced correspondingly. Section~\ref{sec:methods}
describes a third: carve it out of the same control total by vehicle-kilometre share, so the
two categories partition one quantity of diesel.

\subsection{The parallel-table problem}

Chapter~3 of Volume~2 publishes separate default tables for road, rail and off-road mobile
sources, and the same diesel appears in all of them with identical \COtwo{} and sharply
different \NtwoO{}. The distinction reflects real combustion differences---catalyst-equipped
road engines against uncontrolled off-road ones---but the tables are easy to conflate because
the fuel is bought from the same pump.

The national manuscript documents four instances of this trap across the inventory
\citep{ahmad2026national}. Section~\ref{sec:factors} treats the one that arises here, and
reports the diagnostic the engine emits on every run so the consequence is a number rather
than a caution.
